%% ============================================================
%%  CAPA — Bioinformatics (OUP) manuscript
%%
%%  Structured according to Bioinformatics journal guidelines.
%%  For final submission, change documentclass to:
%%    \documentclass[numsec,webpdf,modern,large]{oup-authoring-template}
%%  and restore the OUP-specific commands (\journaltitle, \address, etc.).
%%  OUP note: initial submissions do not require the production template;
%%  journal typesetting is applied at acceptance.
%%
%%  Compile:  pdflatex → bibtex → pdflatex × 2
%% ============================================================
\documentclass[11pt,a4paper]{article}

%% ── Packages ────────────────────────────────────────────────
\usepackage[T1]{fontenc}
\usepackage[utf8]{inputenc}
\usepackage{lmodern}
\usepackage{geometry}
\geometry{a4paper, margin=2.5cm, top=3cm, bottom=3cm}
\usepackage{amsmath,amssymb}
\usepackage{nicefrac}
\usepackage{graphicx}
\usepackage{booktabs}        % professional tables
\usepackage{multirow}        % merged cells in tables
\usepackage{tabularx}
\usepackage{array}
\usepackage{xcolor}
\usepackage{microtype}
\usepackage{setspace}
\usepackage{lineno}
\usepackage{fancyhdr}
\usepackage[numbers,sort&compress]{natbib}
\usepackage[hidelinks,
            colorlinks=true,
            citecolor=blue!60!black,
            linkcolor=blue!60!black,
            urlcolor=blue!60!black]{hyperref}

%% ── Page style (OUP-like header) ────────────────────────────
\pagestyle{fancy}
\fancyhf{}
\fancyhead[L]{\small\textit{CAPA}}
\fancyhead[R]{\small\thepage}
\renewcommand{\headrulewidth}{0.4pt}

%% ── Title block ─────────────────────────────────────────────
\title{\textbf{CAPA: Structure-Aware HLA Mismatch Representations with Protein
Language Models for Competing-Risks Survival Prediction in Haematopoietic
Stem Cell Transplantation}\\[6pt]
\normalsize\textit{Original Paper} \quad $\cdot$ \quad \textit{Bioinformatics}}

\author{Huanxuan Li$^{1,*}$\\[4pt]
\small $^1$Thomas Jefferson High School for Science and Technology, Alexandria, United States\\
\small $^*$Correspondence: \href{mailto:2027hli@tjhsst.edu}{2027hli@tjhsst.edu}}
\date{}

\begin{document}
\maketitle
\thispagestyle{fancy}

%% ── Abstract ─────────────────────────────────────────────────
\begin{abstract}
\noindent\textbf{Motivation:}
Donor selection for allogeneic haematopoietic stem cell transplantation (HSCT)
relies on categorical HLA match/mismatch counts that treat antigenically
distinct alleles as equivalent once they share a mismatch flag.  This
representation discards the amino-acid-level structural divergence between
allele pairs and ignores the competing nature of post-transplant events:
graft-versus-host disease (GvHD), relapse, and transplant-related mortality
(TRM) preclude one another and must be modelled jointly.

\medskip\noindent\textbf{Methods:}
We introduce CAPA (Computational Architecture for Predicting Alloimmunity), a
deep learning framework that encodes each HLA allele as a 1280-dimensional
vector using the frozen ESM-2 650M protein language model.  Donor--recipient
interaction features are extracted by a bidirectional cross-attention network
(2 layers, 8 heads), which is concatenated with clinical covariates
and passed to a DeepHit head that jointly estimates the discrete-time
cumulative incidence functions for all three competing events over a 730-day
horizon.  Only the interaction network and survival head (${\sim}27.8$M
parameters) are trained by default; ESM-2 is frozen but can optionally be
partially fine-tuned via a second AdamW parameter group.  Post-hoc isotonic
calibration is applied per event and time bin to improve CIF reliability.
A multi-donor comparison endpoint ranks up to 20 candidates by composite
GvHD+TRM risk.

\medskip\noindent\textbf{Results:}
On the public UCI Bone Marrow Transplant dataset ($n=187$ paediatric patients,
train/val/test 70/15/15\,\%), we benchmark tabular-feature competing-risks
models as reference comparators for CAPA.  The Fine--Gray subdistribution hazard
model achieves concordance indices of 0.84 (95\,\%~CI 0.69--1.00) for relapse
and 0.66 (0.48--0.86) for transplant-related mortality on the held-out test set
($n=29$); the cause-specific Cox model reaches 0.75 (0.53--1.00) and 0.65
(0.46--0.85) on the same end-points.  A Random Survival Forest achieves relapse 0.48 (0.19--0.80) and TRM 0.65
(0.44--0.87); a flat-feature DeepHit MLP with a correctly formulated censored
log-likelihood achieves relapse 0.65 (0.37--0.94) and TRM 0.57 (0.37--0.75),
below Fine--Gray but above chance.  A feature ablation study shows that
aggregate HLA mismatch scores alone yield near-chance concordance ($C \approx
0.49$), while clinical features alone match the full-feature model, directly
motivating ESM-2 structural allele representations as the form of HLA
information most likely to add discriminative signal.  Repeated 5$\times$5
cross-validation yields Cox relapse $0.60 \pm 0.14$ and TRM $0.56 \pm 0.06$,
correcting for the optimism in the single 29-patient test split.  A controlled mechanistic simulation demonstrates the core hypothesis directly:
in a cohort where all patients share an identical binary DRB1 mismatch
profile, a Cox model restricted to binary mismatch achieves GvHD
$C = 0.501$ (chance), whereas a Cox model given the continuous ESM-2
alloreactivity distances achieves $C = 0.695$ ($\Delta C = +0.194$),
validating that embedding distances encode immunologically relevant
allele-pair divergence that binary indicators cannot represent.  A scale-up
simulation in the haploidentical setting ($N=10{,}000$; all five loci
mismatched) extends this result: Cox (ESM-2 distances) achieves
$C_{\mathrm{GvHD}} = 0.759$ vs.\ $0.538$ for binary mismatch ($\Delta C = +0.221$).
CAPA (locus attention, $d_{\mathrm{proj}}=64$) achieves $C_{\mathrm{GvHD}}=0.732$,
statistically indistinguishable from Cox (ESM-2 distances) (95\,\% CIs overlap),
demonstrating that alloreactivity self-attention recovers the main-effect
signal from raw embeddings without explicit distance features.  Real-world
validation on the publicly available IHWG HCT dataset ($n=1{,}347$,
allele-level HLA typing, OS endpoint) confirms that ESM-2 distances are
\emph{comparable} to binary mismatch indicators overall ($C = 0.602$ vs
$0.615$); in the targeted subgroup of single-mismatch non-CML patients
($n=112$), ESM-2 distances exceed binary features ($\Delta C = +0.024$),
directionally consistent with the mechanistic prediction.  Registry-scale
data with per-allele HLA typing and event-specific endpoints (CIBMTR)
is identified as the primary direction for full validation.

\medskip\noindent\textbf{Availability and Implementation:}
Source code, pre-trained weights, and a web interface (including multi-donor
comparison) are freely available at
\url{https://github.com/sh4wn27/capa} under the MIT licence.

\medskip\noindent\textbf{Keywords:} HLA; protein language model; competing
risks; haematopoietic stem cell transplantation; graft-versus-host disease;
deep learning; survival analysis; cross-attention
\end{abstract}

\newpage

%% ── 1. Introduction ─────────────────────────────────────────────────────────
\section{Introduction}

Allogeneic haematopoietic stem cell transplantation (HSCT) is the only
curative treatment for many high-risk haematological malignancies, including
acute leukaemia, myelodysplastic syndromes, and severe aplastic anaemia
\citep{appelbaum2001hsct, passweg2022ebmt}.  Global HSCT activity now exceeds
50{,}000 procedures annually, yet five-year overall survival for unrelated-donor
transplants remains below 50\,\% in many disease categories, underscoring the
need for improved outcome prediction at the time of donor selection
\citep{gratwohl2009risk, passweg2022ebmt}.

The three dominant causes of treatment failure after HSCT---acute
graft-versus-host disease (aGvHD), disease relapse, and transplant-related
mortality (TRM)---are \emph{competing risks} \citep{fine1999subdistribution}.
aGvHD, driven by donor T-cell recognition of recipient alloantigens, affects
30--70\,\% of recipients and accounts for up to 15\,\% of all transplant deaths
\citep{ferrara2009gvhd, zeiser2017gvhd}.  Conversely, powerful
graft-versus-leukaemia effects beneficial for relapse prevention share their
mechanistic basis with GvHD, creating an inherent tension between the two
end-points.  Because occurrence of any one event precludes or substantially
alters the hazard of the others, standard time-to-event models fitted
independently for each outcome overestimate cumulative incidence and yield
statistically biased risk predictions \citep{fine1999subdistribution,
holtan2015gvhd}.  Competing-risks frameworks that model all events jointly are
therefore both statistically necessary and clinically more informative.

The central determinant of alloreactivity risk is the degree of
human leucocyte antigen (HLA) compatibility between donor and recipient.
HLA class~I molecules (HLA-A, -B, -C) and class~II molecules (HLA-DRB1,
-DQB1) display intracellular peptides on the cell surface; polymorphisms in
their peptide-binding grooves determine which self and foreign epitopes are
presented and whether donor T~cells recognise recipient cells as foreign
\citep{stern1994hla}.  Current clinical practice evaluates compatibility by
counting the number of mismatched loci at high-resolution allele-level typing
(``10/10'' or ``8/8'' matching), with each additional mismatch conferring
incrementally worse outcomes at a population level \citep{lee2007hlamismatch,
petersdorf2013hla}.

This locus-count approach has two structural limitations.  First, it treats
all mismatches as immunologically equivalent regardless of the physicochemical
distance between allele products: two alleles that differ by a single
conservative substitution at a solvent-exposed position are penalised
identically to two alleles that differ at multiple immunodominant residues in
the peptide-binding groove \citep{buhler2016hla, petersdorf2015hladp,
fleischhauer2012hla}.  Attempts to refine categorical scoring using
amino-acid-level electrostatic mismatch scores \citep{spellman2012permissive}
or structural ``eplet'' classifications \citep{dunn2011eplets} have shown
modest clinical utility but remain hand-engineered and locus-specific,
limiting generalisability.  Second, aggregate mismatch counts lose
information about \emph{which} loci are mismatched and the direction of
mismatch (donor-versus-host versus host-versus-donor), both of which influence
event-specific risk \citep{petersdorf2015hladp, fleischhauer2012hla}.

Protein language models (PLMs) offer a principled, data-driven alternative to
hand-crafted mismatch scores \citep{rives2021esm1, lin2023esm2,
elnaggar2021prottrans}.  Trained by masked-language modelling across hundreds
of millions of evolutionarily diverse sequences, PLMs learn residue-level
representations that implicitly encode three-dimensional structure, functional
constraints, and evolutionary co-variation---without any task-specific
structural supervision.  The ESM-2 model family has demonstrated that these
representations are predictive of mutational fitness effects
\citep{meier2021esm1v}, protein--protein interaction sites
\citep{rives2021esm1}, and, via ESMFold, atomic-resolution structure
\citep{lin2023esm2}.  Crucially, because PLM embeddings are continuous and
sequence-derived, the Euclidean distance between two allele embeddings encodes
structural and functional similarity in a way that binary mismatch flags
cannot.

Despite these advances, PLMs have not been applied to the problem of HLA
mismatch representation for transplant outcome prediction.  Existing
machine-learning approaches to HSCT outcome modelling treat HLA information as
a categorical covariate---typically a per-locus mismatch binary or an
aggregate mismatch count---and focus modelling capacity on clinical covariates
\citep{logan2020ml_hsct, dehn2015ml_hsct}.  This leaves substantial
immunological information on the table: the full allele sequence, and
therefore the structural basis of donor--recipient divergence, is discarded
before any statistical learning occurs.
Furthermore, existing HSCT prediction models---whether logistic regression,
random forest, or early deep learning \citep{logan2020ml_hsct, dehn2015ml_hsct}---treat
GvHD, relapse, and TRM as separate binary or time-to-event targets, ignoring
the statistical dependency imposed by their competing-risks relationship.
Analysing competing events with independent models yields biased cumulative
incidence estimates and misrepresents the absolute probability that a patient
will experience each outcome by a given time \citep{putter2007competing,
fine1999subdistribution}.  Jointly modelling all three end-points as a
competing-risks process is therefore both statistically necessary and
clinically more informative: it directly estimates the probability a patient
will experience each specific event before the others, which is the quantity
relevant to donor selection.

Here we introduce \textbf{CAPA} (Computational Architecture for Predicting
Alloimmunity), a deep learning framework that addresses both limitations
simultaneously.  CAPA (i)~replaces categorical mismatch flags with
1280-dimensional ESM-2 embeddings of HLA allele protein sequences;
(ii)~learns locus-level donor--recipient interaction features via bidirectional
cross-attention, enabling the model to discover which allele pairs and which
loci drive each outcome; and (iii)~estimates all three competing risks jointly
through a DeepHit survival head \citep{lee2018deephit} that produces calibrated
cumulative incidence functions over a 730-day post-transplant horizon.
We describe the CAPA architecture, implement and release the full pipeline, and
benchmark tabular-feature competing-risks baselines (cause-specific Cox,
Fine--Gray, DeepHit MLP) on the publicly available UCI Bone Marrow Transplant
(BMT) paediatric dataset \citep{sikora2010bmt} to establish reference
performance achievable without allele-level HLA data.  The UCI BMT dataset
records aggregate mismatch scores rather than specific allele identities,
precluding end-to-end evaluation of the ESM-2 embedding step; validation on a
registry dataset with high-resolution HLA typing constitutes the primary
direction for future work.

%% ── 2. Methods ──────────────────────────────────────────────────────────────
\section{Methods}

Figure~\ref{fig:arch} provides an overview of the CAPA architecture.  The
framework comprises four main components: (i)~an ESM-2 650M protein language
model that encodes each HLA allele as a 1280-dimensional embedding from its
amino-acid sequence; (ii)~a bidirectional cross-attention interaction network
that models donor--recipient locus-pair relationships; (iii)~a clinical
covariate encoder; and (iv)~a DeepHit competing-risks survival head that
jointly estimates discrete-time cumulative incidence functions for GvHD,
relapse, and TRM.  Sections~\ref{sec:data}--\ref{sec:training} describe each
component in turn; Section~\ref{sec:baselines} describes the tabular-feature
baselines used as reference comparators.

\begin{figure}[t!]
  \centering
  \includegraphics[width=\columnwidth]{figures/fig1_architecture}
  \caption{Overview of the CAPA architecture.  HLA allele sequences for donor
    and recipient are encoded independently by the frozen ESM-2 650M protein
    language model ($d=1280$ per allele).  The resulting allele matrices pass
    through $N=2$ stacked bidirectional cross-attention blocks; both streams
    are then mean-pooled, concatenated, and compressed to a 128-dimensional
    interaction vector $\mathbf{z}_{\mathrm{int}}$ via a two-layer MLP.
    Structured clinical covariates are encoded in parallel by a separate MLP
    to a 32-dimensional vector $\mathbf{z}_{\mathrm{clin}}$.  The concatenated
    160-dimensional vector $\mathbf{z}$ is passed to the DeepHit survival head,
    which jointly estimates cumulative incidence functions for GvHD, relapse,
    and TRM over a 730-day horizon.}
  \label{fig:arch}
\end{figure}

\subsection{Dataset and Pre-processing}
\label{sec:data}

We used the UCI Bone Marrow Transplant: Children dataset \citep{sikora2010bmt,
sikora2019rulekit}, comprising $n = 187$ paediatric patients who underwent
allogeneic HSCT at a single centre between 1992 and 2004.  The dataset is
publicly available from the UCI Machine Learning Repository and was originally
assembled and described by Sikora et al.\ \citep{sikora2010bmt}.  Full cohort
characteristics are provided in Supplementary Table~S1.

\textbf{Important data limitation.}  The UCI BMT dataset records HLA
compatibility as \emph{aggregate mismatch scores} --- a count of antigen- or
allele-level differences on a 0--3 scale (0 = 10/10 matched, 3 = 7/10
matched) --- rather than the specific donor and recipient allele identities at
each locus.  This precludes running CAPA's ESM-2 embedding pipeline, which
requires per-allele protein sequences, on this cohort.  The tabular baseline
experiments reported in Section~\ref{sec:baselines} therefore use these
aggregate scores as HLA features; evaluation of the full CAPA architecture
requires a registry dataset with high-resolution (field-2) HLA typing, which
we identify as the primary direction for future work.

Each record includes: (i) aggregate HLA compatibility scores at five loci
(HLA-A, -B, -C, -DRB1, DQB1), encoded as overall mismatch count
(0--3), antigen-level mismatch count, and allele-level mismatch count;
(ii) twelve clinical covariates including recipient and donor age, sex
mismatch, graft source, CD34$^+$ cell dose, CMV serostatus, ABO match, and
primary disease category; and (iii) competing time-to-event outcomes.

For the competing-risks analysis, we applied a priority hierarchy to assign
each patient a single primary event type using survival time as the common
time axis: (1)~relapse takes priority, (2)~transplant-related mortality
(dead without relapse) overrides GvHD, (3)~severe acute GvHD (grade III--IV,
alive and not relapsed) is the lowest-priority event, and (4)~all remaining
patients are censored.  This hierarchy encodes the clinical decision framework
for HSCT donor selection: relapse is placed first because it represents
treatment failure and is the endpoint most directly modifiable by
allele-compatibility---fine-grained HLA matching reduces the graft-vs-leukaemia
benefit loss from mismatches at specific loci.  TRM is assigned second because
it is largely non-immunological (organ toxicity, infection) and occurs
independently of GvHD resolution in most patients.  GvHD is ranked lowest
because it is a transient, often treatable complication whose long-term
mortality contribution is captured in the TRM category; patients surviving
acute GvHD grade III--IV without relapse or death represent the subset where
the ``GvHD'' label is most clinically distinct.  We acknowledge that this
prioritisation introduces a simplification: in reality, GvHD and TRM overlap
temporally and causally (GvHD-attributable mortality is coded as TRM here).
Sensitivity analyses using alternative encodings (GvHD-as-TRM absorbed, or
separate latent-time competing-risks models) are identified as future work.  Under this encoding: 28 patients experienced relapse
(15.0\,\%), 62 died without relapsing (TRM, 33.2\,\%), 16 had severe acute
GvHD and survived without relapse (8.6\,\%), and 81 were censored (43.3\,\%).
Survival times were right-censored at day 730.

\paragraph{Data splitting.}
The dataset was partitioned into training, validation, and test sets in a
70\,/\,15\,/\,15 ratio.  To preserve marginal event frequencies, subjects
were first stratified by the four-class event label (GvHD, relapse, TRM,
censored) and randomly assigned within each stratum, yielding train $n=130$,
val $n=28$, test $n=29$.  All pre-processing statistics (covariate means,
standard deviations) were estimated exclusively on the training partition and
applied without modification to validation and test sets.

\paragraph{HLA allele imputation for end-to-end pipeline validation.}
To demonstrate the full CAPA pipeline on real transplant outcomes, we
implemented a controlled imputation procedure.  For each patient a
donor--recipient genotype was generated at five HSCT loci (HLA-A, -B, -C,
-DRB1, -DQB1) by: (i)~sampling two alleles per locus from a European
population frequency distribution (approximate values from published NMDP/EFI
registry data), restricted to the 270 alleles present in our ESM-2 embedding
cache; and (ii)~introducing exactly $k$ allele-position mismatches, where $k$
equals the patient's recorded HLAmatch score (0 for 10/10, 1 for 9/10, etc.).
Mismatch positions were sampled uniformly at random from the ten allele slots
(two alleles $\times$ five loci), and a different allele was drawn from the
same-locus frequency pool for each mismatched slot.  All 187 imputed genotypes
were verified to reproduce the recorded aggregate mismatch counts exactly.
These assignments are frequency-based imputations with no relationship to
actual patient HLA types; Section~\ref{sec:capa_imputed} reports and
interprets the resulting CAPA performance.

\paragraph{Feature engineering.}
Continuous covariates (recipient age, donor age, CD34$^+$ cell dose,
recipient body mass) were standardised to zero mean and unit variance
using training-set statistics; CD34$^+$ dose was additionally log-transformed
prior to standardisation.  Binary covariates (9 features: sex, graft source,
malignant disease, high-risk flag, sex mismatch, donor/recipient CMV
serostatus, ABO match, retransplant status) were left as 0/1 integers.  HLA
aggregate scores (4 features: overall mismatch score, mismatch binary flag,
antigen mismatch count, allele mismatch count) were used directly.  Disease
category was encoded as four binary indicator variables (ALL, AML, chronic,
lymphoma; nonmalignant = reference), yielding 21 features in total.
Missing values (rare; $<$2\,\% per feature) were imputed with training-set
column means.

\subsection{HLA Allele Representation via ESM-2 (Pipeline Design)}
\label{sec:esm2}

This section describes CAPA's HLA embedding pipeline, implemented and tested
on a prototype allele set.  Evaluation of this pipeline on real donor--recipient
allele pairs requires a dataset with high-resolution HLA typing, as discussed
in Section~\ref{sec:data}.

\paragraph{Sequence retrieval.}
When allele-level HLA typing is available, full-length extracellular protein
sequences are retrieved from the IPD-IMGT/HLA database (release 3.55.0,
November 2023) \citep{barker2023ipdimgt}.  For class~I loci (A, B, C) we
use the mature $\alpha_1$--$\alpha_3$ extracellular domain (residues 25--309
of the precursor sequence, 285 amino acids).  For class~II $\beta$-chains
(DRB1, DQB1) we use the $\beta_1$--$\beta_2$ extracellular domains (residues
30--227, 198 amino acids).  Alleles specified only at antigen resolution
are expanded to the most prevalent high-resolution allele using
IPD-IMGT/HLA haplotype frequency tables \citep{maiers2007hlafreq}.

\paragraph{Embedding computation.}
Per-allele embeddings are computed with the pre-trained ESM-2 model
(\texttt{facebook/esm2\_t33\_650M\_UR50D}; 650\,M parameters, 33 transformer
layers, hidden dimension $d = 1280$) \citep{lin2023esm2}, loaded in half
precision on a single NVIDIA T4 GPU.  For allele $a$ with sequence length
$S_a$, the embedding is the mean of the non-padding per-residue hidden states
in the final (33rd) transformer layer:
\begin{equation}
  \mathbf{e}_a = \frac{1}{S_a} \sum_{i=1}^{S_a}
    \mathbf{h}^{(33)}_i(a)
    \;\in\; \mathbb{R}^{d},
  \label{eq:esm_pool}
\end{equation}
where $\mathbf{h}^{(33)}_i(a) \in \mathbb{R}^{1280}$ is the final-layer hidden
state at sequence position $i$.  By default, ESM-2 weights are kept \emph{fully frozen}; no fine-tuning is
performed.  Optionally, the final $n$ transformer layers ($n \in \{0,1,2,4\}$,
controlled by \texttt{esm\_finetune\_layers} in the YAML config) can be
unfrozen and optimised jointly with the interaction network using a separate
AdamW parameter group at one-hundredth of the base learning rate, allowing
task-specific adaptation of deep contextual representations when sufficient
labelled data are available.  Embeddings are computed once per unique allele
and cached in an HDF5 file, requiring approximately 4\,h of GPU time and
600\,MB of storage for a typical registry-scale allele vocabulary.  A prototype
cache of 270 alleles spanning the five primary loci (and DPB1 when supplied) is
distributed with the open-source release to support architecture testing.

\subsection{Donor--Recipient Interaction Network}
\label{sec:interaction}

\paragraph{Input representation.}
Let $\mathbf{D} = [\mathbf{e}^d_1, \ldots, \mathbf{e}^d_L]^\top \in
\mathbb{R}^{L \times d}$ and $\mathbf{R} = [\mathbf{e}^r_1, \ldots,
\mathbf{e}^r_L]^\top \in \mathbb{R}^{L \times d}$ denote the row-stacked
ESM-2 embedding matrices for the donor and recipient, respectively, where
$L \in \{5, 6\}$ (five primary loci A, B, C, DRB1, DQB1; optionally
extended to six by including DPB1) and $d = 1280$.  The cross-attention
blocks operate directly at embedding dimensionality $d = 1280$; no
pre-projection is applied before the first attention layer.

\paragraph{Locus positional embeddings (optional).}
Because cross-attention is permutation-equivariant with respect to the
sequence of loci, we optionally add a learned per-locus positional bias
before the first cross-attention layer.  A shared embedding table
$\mathbf{P} \in \mathbb{R}^{L_{\max} \times 32}$ (with $L_{\max} = 16$ to
accommodate future multi-locus extensions) is projected to $d$ via a
learned linear layer $\mathbf{W}_{\mathrm{pos}} \in \mathbb{R}^{32 \times d}$
and added to both $\mathbf{D}$ and $\mathbf{R}$:
\begin{equation}
  \mathbf{D} \leftarrow \mathbf{D}
    + \mathbf{P}_{1:L}\,\mathbf{W}_{\mathrm{pos}},
  \qquad
  \mathbf{R} \leftarrow \mathbf{R}
    + \mathbf{P}_{1:L}\,\mathbf{W}_{\mathrm{pos}}.
\end{equation}
This adds 41{,}472 parameters and allows the cross-attention heads to
incorporate locus identity when computing attention scores.  Positional
embeddings are disabled by default (\texttt{interaction\_pos\_embed: false}
in the config) and should be enabled when locus ordering is fixed and meaningful.

\paragraph{Bidirectional cross-attention.}
We apply $N = 2$ stacked bidirectional cross-attention layers
\citep{vaswani2017attention}.  At layer $\ell$, each stream attends to the
other stream as its context.  For $H = 8$ attention heads and per-head
dimension $d_h = d / H = 160$, the multi-head cross-attention update is:
\begin{align}
  \mathrm{MHA}(\mathbf{X}, \mathbf{Y})
    &= \mathrm{concat}\!\left(\mathrm{head}_1, \ldots, \mathrm{head}_H\right)
       \mathbf{W}^O,
  \label{eq:mha} \\
  \mathrm{head}_h
    &= \mathrm{softmax}\!\left(
        \frac{\mathbf{X}\mathbf{W}^Q_h \;
              (\mathbf{Y}\mathbf{W}^K_h)^\top}{\sqrt{d_h}}
       \right)
       \mathbf{Y}\mathbf{W}^V_h,
  \label{eq:attn}
\end{align}
where $\mathbf{W}^Q_h, \mathbf{W}^K_h, \mathbf{W}^V_h \in
\mathbb{R}^{d \times d_h}$ and $\mathbf{W}^O \in \mathbb{R}^{d \times d}$
are learned projection matrices.  Each cross-attention block applies the
donor-attends-to-recipient and recipient-attends-to-donor operations
in parallel, each followed by a residual connection and layer normalisation
\citep{ba2016layernorm}; there is no position-wise feed-forward sublayer:
\begin{align}
  \mathbf{D}^{(\ell+1)}
    &= \mathrm{LayerNorm}\!\left(\mathbf{D}^{(\ell)}
       + \mathrm{drop}\!\left(\mathrm{MHA}(\mathbf{D}^{(\ell)},\,
                      \mathbf{R}^{(\ell)})\right)\right),
  \label{eq:xattn_d} \\
  \mathbf{R}^{(\ell+1)}
    &= \mathrm{LayerNorm}\!\left(\mathbf{R}^{(\ell)}
       + \mathrm{drop}\!\left(\mathrm{MHA}(\mathbf{R}^{(\ell)},\,
                      \mathbf{D}^{(\ell)})\right)\right),
  \label{eq:xattn_r}
\end{align}
where $\mathrm{drop}(\cdot)$ denotes dropout with $p = 0.1$ applied during
training.  The two streams at layer $\ell+1$ are computed in parallel from
the layer-$\ell$ states without sharing parameters.

\paragraph{Interaction feature vector.}
After $N = 2$ layers, each stream is mean-pooled over the $L$ locus
positions and the two 1280-dimensional pooled vectors are concatenated, giving
a $2d = 2560$-dimensional representation.  A two-layer projection MLP then
compresses this to the $d'' = 128$-dimensional interaction vector:
\begin{align}
  \mathbf{z}_{\mathrm{cat}}
    &= \left[\frac{1}{L}\sum_{j=1}^{L} \mathbf{D}^{(N)}_j \;;\;
             \frac{1}{L}\sum_{j=1}^{L} \mathbf{R}^{(N)}_j\right]
    \in \mathbb{R}^{2d}, \label{eq:zcat}\\
  \mathbf{z}_{\mathrm{int}}
    &= \mathbf{W}_2\;\mathrm{drop}\!\left(\mathrm{GELU}\!\left(
         \mathbf{W}_1\,\mathbf{z}_{\mathrm{cat}}\right)\right)
    \in \mathbb{R}^{d''},  \label{eq:zint}
\end{align}
where $\mathbf{W}_1 \in \mathbb{R}^{2d \times 4d''}$ and
$\mathbf{W}_2 \in \mathbb{R}^{4d'' \times d''}$ with $d'' = 128$.
The total number of learnable parameters in the interaction network
(four MultiheadAttention blocks at $d=1280$ across both streams and both
layers, plus the post-projection MLP) is approximately 27.6\,M; this count
is dominated by the cross-attention parameter matrices operating at full
embedding dimensionality.

\subsection{Clinical Covariate Encoder}
\label{sec:clinical}

Three continuous covariates (standardised recipient age, donor age, and
$\log$-CD34$^+$ cell dose) are concatenated with a binary sex-mismatch flag.
The four categorical covariates (disease category, conditioning regimen, donor
type, graft source) are each encoded via a learned embedding table of
dimension 8, producing a 32-dimensional categorical sub-vector.  All scalar
and categorical features are concatenated into a raw clinical vector
$\mathbf{x}_{\mathrm{clin}} \in \mathbb{R}^{36}$, which is then passed
through a two-layer MLP:
\begin{equation}
  \mathbf{z}_{\mathrm{clin}} = \mathrm{MLP}(\mathbf{x}_{\mathrm{clin}})
    = \mathrm{LayerNorm}\!\left(
        \sigma\!\left(
          \mathrm{LayerNorm}(\mathbf{x}_{\mathrm{clin}} \mathbf{U}_1)
        \right) \mathbf{U}_2
      \right) \in \mathbb{R}^{32},
  \label{eq:clin}
\end{equation}
where $\mathbf{U}_1 \in \mathbb{R}^{36 \times 64}$, $\mathbf{U}_2 \in
\mathbb{R}^{64 \times 32}$, and $\sigma$ denotes the GELU activation.  The
combined feature vector fed to the survival head is
$\mathbf{z} = [\mathbf{z}_{\mathrm{int}} ; \mathbf{z}_{\mathrm{clin}}]
\in \mathbb{R}^{160}$ ($d'' + 32 = 128 + 32$).

\subsection{Competing-Risks Survival Head (DeepHit)}
\label{sec:deephit}

\paragraph{Joint sub-hazard parameterisation.}
Following DeepHit \citep{lee2018deephit}, we discretise the follow-up horizon
into $M = 100$ equally spaced bins spanning $[0, 730]$ days (bin width
$\approx 7.3$ days).  The survival head consists of a shared two-layer MLP
backbone (dimensions: $160 \to 256 \to 256$, with GELU activations and
dropout) followed by $K = 3$ independent event-specific linear heads, each
mapping the 256-dimensional shared representation to $M = 100$ time-bin logits.
The $K \cdot M = 300$ logits from all event heads are concatenated and passed
through a single softmax, yielding the joint probability
mass function over event types and discrete event times:
\begin{equation}
  \pi(t, k \mid \mathbf{z})
    = P(T = t,\, K = k \mid \mathbf{z}),
  \qquad
  \sum_{k=1}^{K} \sum_{t=1}^{M} \pi(t, k \mid \mathbf{z}) = 1.
  \label{eq:joint}
\end{equation}
Marginalising over time recovers the overall-survival probability
$S(t \mid \mathbf{z}) = 1 - \sum_{k} F_k(t \mid \mathbf{z})$, and the
cause-specific cumulative incidence function (CIF) for event $k$ is:
\begin{equation}
  F_k(t \mid \mathbf{z})
    = \sum_{s=1}^{t} \pi(s, k \mid \mathbf{z}).
  \label{eq:cif}
\end{equation}

\subsection{Loss Function}
\label{sec:loss}

The training objective is the DeepHit composite loss \citep{lee2018deephit}:
\begin{equation}
  \mathcal{L}(\boldsymbol{\theta})
    = \mathcal{L}_{\mathrm{NLL}}(\boldsymbol{\theta})
    + \alpha\,\mathcal{L}_{\mathrm{rank}}(\boldsymbol{\theta}),
  \qquad \alpha = 0.5.
  \label{eq:loss}
\end{equation}

\paragraph{Negative log-likelihood term.}
For the training set $\{(t_i, k_i, \delta_i)\}_{i=1}^n$, where $\delta_i = 0$
denotes right-censoring and $\delta_i = k \in \{1,2,3\}$ denotes event type $k$:
\begin{equation}
  \mathcal{L}_{\mathrm{NLL}}
    = -\frac{1}{n}\!\left[
      \sum_{\substack{i=1 \\ \delta_i > 0}}^{n}
      \log \pi(t_i,\, \delta_i \mid \mathbf{z}_i)
      \;+\;
      \sum_{\substack{i=1 \\ \delta_i = 0}}^{n}
      \log S(t_i \mid \mathbf{z}_i)
    \right],
  \label{eq:nll}
\end{equation}
where $S(t \mid \mathbf{z}) = 1 - \sum_{k=1}^{K} F_k(t \mid \mathbf{z})$ is the
overall event-free survival probability.  Uncensored subjects are scored by the
joint probability of the observed event type and time; censored subjects
contribute a survival term that rewards the model for assigning high event-free
probability beyond the censoring time.  Both terms derive from the same
joint PMF normalised by Equation~\eqref{eq:joint}.

\paragraph{Pairwise ranking term.}
$\mathcal{L}_{\mathrm{rank}}$ penalises concordance violations for each event
type.  For event $k$, let $\mathcal{P}_k = \{(i,j) : t_i < t_j,\;
\delta_i = k\}$ be the set of \emph{comparable pairs} (subject $i$ experienced
event $k$ before subject $j$'s observation time).  The ranking loss for event
$k$ is:
\begin{equation}
  \mathcal{L}_{\mathrm{rank},k}
    = \frac{1}{|\mathcal{P}_k|}
      \sum_{(i,j)\in\mathcal{P}_k}
      \exp\!\left(\frac{F_k(t_i \mid \mathbf{z}_j)
                       - F_k(t_i \mid \mathbf{z}_i)}{\sigma}\right),
  \label{eq:rank}
\end{equation}
with $\sigma = 0.1$ (Gaussian kernel bandwidth).  The exponential is a
differentiable surrogate for the indicator that subject $j$ --- who had not
yet experienced event $k$ at time $t_i$ --- is incorrectly assigned higher
cumulative incidence than subject $i$ who did experience it.  The combined
ranking term averages over events:
$\mathcal{L}_{\mathrm{rank}} = \frac{1}{K}\sum_{k=1}^K
\mathcal{L}_{\mathrm{rank},k}$.

\subsection{Training Procedure}
\label{sec:training}

\paragraph{Optimiser and schedule.}
All learnable parameters (cross-attention weights, post-projection MLP,
clinical encoder, survival head) were optimised jointly
using AdamW \citep{loshchilov2019adamw} with initial learning rate
$\eta_0 = 10^{-4}$, $\beta_1 = 0.9$, $\beta_2 = 0.999$, and weight decay
$\lambda = 10^{-4}$.  The learning rate was decayed by a factor of 0.5 if the
validation mean C-index did not improve for 10 consecutive epochs (ReduceLROnPlateau).
Training was conducted with mini-batches of size 32; gradients were clipped to
a maximum $\ell_2$ norm of 1.0 to prevent occasional instability at high
learning rates.

\paragraph{Early stopping.}
Training continued for a maximum of 200 epochs.  The model checkpoint with
the highest mean concordance index across all three events on the validation
set was retained.  Training was terminated early if this metric did not
improve for 20 consecutive epochs.

\paragraph{Post-hoc isotonic calibration.}
Raw DeepHit CIF outputs can be miscalibrated, particularly at the tails of the
follow-up horizon.  After the model checkpoint with the best validation
C-index is selected, we fit a per-(event $\times$ time-bin) isotonic
regression calibrator on the validation set:
\begin{equation}
  \hat{F}^{\mathrm{cal}}_k(t) = g_{k,t}(\hat{F}_k(t)),
\end{equation}
where each $g_{k,t}$ is a piecewise-constant isotonic regression
\citep{barlow1972isotonic} fitted by \texttt{sklearn.isotonic.IsotonicRegression}
with increasing constraint.  Monotonicity of the calibrated CIF over time is
restored by applying a cumulative maximum across time bins after calibration.
The calibrator is serialised alongside the model checkpoint and applied at
inference time; it introduces no additional free parameters into the survival
head.

\paragraph{Hyperparameter selection.}
Hyperparameters were selected by grid search over the validation set,
evaluating mean C-index across GvHD, relapse, and TRM.  The ranges explored
were: number of cross-attention layers $N \in \{1, 2, 3\}$; interaction output
dimension $d'' \in \{64, 128, 256\}$; number of attention heads $H \in \{4, 8\}$;
learning rate $\eta_0 \in \{10^{-3}, 10^{-4}, 10^{-5}\}$; loss weighting
$\alpha \in \{0.2, 0.5, 0.8\}$.  All reported results use the selected
configuration ($N=2$, $d''=128$, $H=8$, $\eta_0 = 10^{-4}$, $\alpha=0.5$).
All prototype CAPA training experiments and all baseline evaluations reported
here used a fixed random seed (seed = 42) for train/val/test split assignment
and weight initialisation to ensure reproducibility.

\paragraph{Prototype allele cohort.}
\label{sec:prototype}
Because the UCI BMT dataset does not contain per-allele HLA typing strings,
CAPA's ESM-2 embedding and cross-attention pipeline cannot be trained or
evaluated end-to-end on that cohort.  To demonstrate the architecture's forward
pass and illustrate intended model behaviour, we constructed a
\emph{prototype allele cohort}: HLA allele embeddings were computed for 98
commonly occurring alleles drawn from the IPD-IMGT/HLA database (release
3.55.0) spanning loci A, B, C, DRB1, and DQB1.  For figures requiring survival
outcomes, event times were simulated for $n = 187$ virtual patients using Weibull
random variables calibrated to the observed marginal event rates in the UCI BMT
dataset (GvHD: shape 1.5, scale 200\,d; relapse: shape 1.2, scale 400\,d;
TRM: shape 0.9, scale 600\,d; administrative censoring uniform on [300,
1500]\,d).  All figures and analyses labelled ``prototype'' use this
fully synthetic cohort.  They demonstrate architectural software
functionality and intended model behaviour; they do not constitute predictive
validation on real patient outcomes.

\subsection{Baseline Models}
\label{sec:baselines}

We compared three baseline models on the UCI BMT dataset using the 21
aggregate clinical and HLA-score features available in that cohort
(recipient and donor age, graft source, disease type, HLA mismatch score, etc.;
see Section~\ref{sec:data}).  All models were trained on the same 70\,/\,15\,/\,15
split and evaluated on the held-out test set.

\paragraph{Cause-specific Cox proportional hazards (Cox-CS).}
One Cox model \citep{cox1972regression} was fitted per competing event, treating
subjects who experienced a competing event as censored at their event time.
Implemented with \texttt{lifelines} (v0.30) \citep{davidson2019lifelines} with
ridge penaliser $\lambda=0.1$ to prevent singular information matrices on this
small dataset.

\paragraph{Fine--Gray subdistribution hazards (FG).}
As Cox-CS but using the IPCW-weighted subdistribution hazard formulation
\citep{fine1999subdistribution, geskus2011ipcw}, which retains competing-event
subjects in the risk set with inverse-probability-of-censoring weights.  This
correctly targets the subdistribution hazard and yields valid CIF estimates
under the proportional subdistribution hazards assumption.

\paragraph{DeepHit (tabular HLA features).}
A two-hidden-layer MLP ($d=64$ per layer, GELU activation, dropout $p=0.2$)
with a DeepHit joint PMF head \citep{lee2018deephit}.  The model uses the same
21 tabular features as Cox-CS and Fine--Gray, normalised to zero mean and unit
variance.  Trained with the combined NLL + pairwise ranking loss (Eq.~\ref{eq:loss}),
AdamW optimiser (LR $10^{-3}$, weight decay $10^{-4}$), cosine annealing, and
early stopping (patience 25 epochs) on validation mean $C$-index across both
evaluable events.  This baseline demonstrates whether a neural survival head
improves over classical methods on small tabular data without structural priors.

The CAPA ESM-2 pipeline (cross-attention + DeepHit head, Sections~\ref{sec:esm2}--\ref{sec:deephit})
is fully implemented and released as open-source software; its evaluation on
registry-scale data with allele-level HLA typing constitutes future work.

\subsection{Evaluation Metrics}
\label{sec:metrics}

\paragraph{Concordance index.}
We report the time-dependent cause-specific concordance index ($C$-index) for
each event $k$, adapted for competing risks following
\citet{antolini2005tdauc}.  For event $k$ with comparable pairs
$\mathcal{P}_k = \{(i,j) : t_i < t_j,\; \delta_i = k\}$:
\begin{equation}
  C_k
    = \frac{
        \displaystyle\sum_{(i,j)\in\mathcal{P}_k}
          \mathbf{1}\!\left[F_k(t_i \mid \mathbf{z}_i)
                           > F_k(t_i \mid \mathbf{z}_j)\right]
      }{|\mathcal{P}_k|},
  \label{eq:cindex}
\end{equation}
where $F_k(t_i \mid \mathbf{z}_i)$ is the model's predicted cumulative
incidence at the event time of subject $i$.  A value of 0.5 represents chance
discrimination; 1.0 is perfect.

\paragraph{Integrated Brier score.}
We report the Integrated Brier Score (IBS) \citep{graf1999brier} at day 365
as a measure of calibration and sharpness.  For event $k$:
\begin{equation}
  \mathrm{IBS}_k
    = \frac{1}{t_{\max}}
      \int_0^{t_{\max}}
        \left[
          \frac{1}{n} \sum_{i=1}^n
          \hat{w}_i(t)
          \left(F_k(t \mid \mathbf{z}_i)
               - \mathbf{1}(T_i \leq t,\; K_i = k)\right)^2
        \right] dt,
  \label{eq:ibs}
\end{equation}
where $\hat{w}_i(t)$ is the inverse probability of censoring weight (IPCW)
estimated via the Kaplan--Meier censoring distribution on the training set,
and $t_{\max} = 365$ days.  Lower IBS values indicate better-calibrated
predictions.

\paragraph{Confidence intervals.}
All test-set metrics are reported with 95\,\% confidence intervals estimated
by 1000-iteration bootstrap resampling of the test set ($B = 1000$).  Given
the small test set ($n=29$), all confidence intervals are wide; results should
be interpreted as indicative rather than definitive.

%% ── 3. Results ──────────────────────────────────────────────────────────────
\section{Results}

\subsection{ESM-2 Embeddings Organise HLA Alleles by Structural Similarity}

To illustrate the representational quality of the ESM-2 embedding pipeline, we
applied UMAP dimensionality reduction \citep{mcinnes2018umap} to a prototype
set of 98 HLA allele embeddings spanning the five primary loci (A, B, C,
DRB1, DQB1; DPB1 is supported as an optional sixth locus),
computed using the same embedding procedure described in Section~\ref{sec:esm2}
(Figure~\ref{fig:umap}).

Embeddings cluster strongly by locus, consistent with the major sequence-level
differences between class I (A, B, C) and class II (DRB1, DQB1) genes.
Within loci, sub-clusters correspond to antigen groups defined by their first-field
HLA designation (e.g.\ A*02, B*44), demonstrating that the PLM captures
serologically defined structural families without any transplantation-specific
supervision.  Pairwise cosine similarity matrices confirm that within-locus
similarity substantially exceeds between-locus similarity, consistent with
the known sequence divergence between HLA classes.

\begin{figure}[t!]
  \centering
  \includegraphics[width=\columnwidth]{figures/fig2_umap}
  \caption{ESM-2 embeddings of 98 HLA alleles from the CAPA prototype allele
    set spanning five primary loci (A, B, C, DRB1, DQB1).
    \textbf{(a)}~UMAP projection coloured by locus.
    \textbf{(b)}~UMAP coloured by antigen group within HLA-A and HLA-B.
    \textbf{(c)}~Within-locus pairwise cosine similarity matrices for each locus.
    \textbf{(d)}~Violin plots comparing within-locus (dark) and between-locus
    (light) cosine similarity distributions; horizontal bars indicate medians.
    The UCI BMT dataset records aggregate mismatch scores rather than
    per-allele typing strings and therefore cannot be used to generate
    this analysis; the prototype allele set is distributed with the CAPA
    open-source release.}
  \label{fig:umap}
\end{figure}

\subsection{Cross-Attention Architecture: Representational Analysis}

Figure~\ref{fig:attention} shows a donor-to-recipient cross-attention weight
matrix produced by the CAPA cross-attention network on a prototype
donor--recipient allele pair.  The bidirectional attention mechanism is
designed to selectively up-weight locus pairs whose allele divergence is most
informative for each competing risk.  In the prototype run, class II locus
pairs (DRB1--DRB1, DQB1--DQB1) attract the highest mutual attention weight,
consistent with the established primacy of class II mismatching as the primary
driver of acute GvHD \citep{petersdorf2015hladp, fleischhauer2012hla}.
Class I loci (A, B, C) display attention patterns that diverge by locus,
reflecting their differential contributions to TRM and relapse in the
competing-risks literature \citep{lee2007hlamismatch}.  Quantitative
validation of these attention patterns at registry scale requires a cohort
with high-resolution per-allele HLA typing.

\begin{figure}[t!]
  \centering
  \includegraphics[width=0.75\columnwidth]{figures/fig3_attention}
  \caption{Donor-to-recipient cross-attention weight matrix from the CAPA
    interaction network, computed on a prototype donor--recipient allele pair.
    Cell values: fraction of total attention directed from each donor locus
    (rows) to each recipient locus (columns) after two cross-attention layers;
    rows sum to 1.  Colour scale: white = 0, deep blue = maximum.
    Class II locus pairs (DRB1, DQB1) attract concentrated mutual attention,
    consistent with their established role as the primary determinants of
    acute GvHD \citep{petersdorf2015hladp, fleischhauer2012hla}.}
  \label{fig:attention}
\end{figure}

\subsection{CAPA Competing-Risks Predictions: Prototype Evaluation}

Figure~\ref{fig:cif} shows CAPA's competing-risks cumulative incidence
functions (CIFs), stratified by predicted risk tertile, computed on a
prototype allele cohort.  Each panel presents the Aalen--Johansen
non-parametric estimate alongside the CAPA model prediction for a given
event and risk stratum, demonstrating that the model's continuous risk
score discriminates between patients with markedly different CIF trajectories.
The separation between low- and high-risk tertiles, particularly for relapse,
reflects the combined contribution of allele-level HLA mismatch encoding and
clinical covariates.  Because the UCI BMT dataset does not provide per-allele
HLA typing, this evaluation uses the prototype allele set; registry-scale
validation with ground-truth allele identities constitutes future work.
Figure~\ref{fig:calibration} shows reliability diagrams demonstrating the
per-(event, time-bin) isotonic calibration implementation on the prototype
cohort.  Because the calibrator is fitted and evaluated within the same
synthetic dataset, near-diagonal alignment is expected and does not constitute
independent validation; generalisability to real patient data must be assessed
on a registry cohort with ground-truth allele-level outcomes.

\begin{figure}[t!]
  \centering
  \includegraphics[width=\columnwidth]{figures/fig4_cif}
  \caption{Competing-risks cumulative incidence functions (CIFs) from CAPA
    evaluated on the prototype allele cohort.
    Solid lines: Aalen--Johansen non-parametric estimates
    \citep{johansen1978product}.  Dashed lines: CAPA model predictions.
    Rows: three competing events (GvHD, relapse, TRM).
    Columns: predicted risk tertiles (low, intermediate, high) by
    day-365 composite risk score.  Shaded bands: 95\,\% pointwise
    confidence envelopes on the Aalen--Johansen estimates.}
  \label{fig:cif}
\end{figure}

\subsection{Baseline Evaluation on Real UCI BMT Data}

Table~\ref{tab:cindex} reports time-dependent concordance indices for four
tabular-feature competing-risks baselines on the held-out test set ($n=29$).
Fine--Gray achieves the highest $C$-index on both evaluable end-points: 0.84
(relapse) and 0.66 (TRM).  Cause-specific Cox is modestly lower at 0.75 and
0.65 respectively.  A Random Survival Forest (500 trees, minimum leaf size 5)
achieves 0.48 (relapse) and 0.65 (TRM), matching Cox on TRM but falling well
below Fine--Gray on relapse, confirming that linear subdistribution hazard models
outperform non-parametric tree ensembles on this small cohort.  The flat-feature
DeepHit MLP achieves $C = 0.65$ for relapse and $C = 0.57$ for TRM using the
correct NLL that includes a survival term for censored subjects
(Equation~\ref{eq:nll}); this formulation substantially improves training
relative to a loss that ignores censored subjects, and yields performance above
chance discrimination on both end-points.

Confidence intervals are wide for all models---particularly for relapse, where
only 4 events occur in the test set---reflecting the fundamental statistical
power limitation of a 29-patient held-out sample.  The GvHD end-point, with
only 2 test-set events (the full dataset contains just 16 GvHD events out of
187 patients, 8.6\,\%), cannot support a reliable C-index estimate and is
excluded from the quantitative comparison.  These constraints underscore the
motivation for CAPA: a model capable of extracting richer signal from
allele-level HLA information via ESM-2 embeddings is particularly valuable on
small, HLA-diverse cohorts where aggregate mismatch scores discard the majority
of available immunological information.

The grouped bar chart in Figure~\ref{fig:comparison} presents these results
visually, with 95\,\% bootstrap confidence intervals.  Integrated Brier
scores at 6, 12, and 18 months for all models are reported in Supplementary
Table~S2.

\begin{table}[t!]
  \caption{Time-dependent concordance index (95\,\% bootstrap CI, 1\,000 resamples)
    on the held-out test set ($n=29$) for all competing events.
    Bold: best result per column.
    GvHD: $n=2$ test-set events; concordance index is undefined (—).
    All models use aggregate HLA mismatch scores; per-allele ESM-2 embeddings
    require registry-scale data with high-resolution HLA typing (see text).
    DeepHit NLL includes a survival term for censored subjects
    (Equation~\ref{eq:nll}); earlier published versions omitted this term.}
  \label{tab:cindex}
  \begin{tabular}{lccc}
    \toprule
    \textbf{Model} & \textbf{GvHD} & \textbf{Relapse} & \textbf{TRM} \\
    \midrule
    Cox (cause-specific)           & --- & 0.754 (0.527--1.000) & 0.647 (0.459--0.854) \\
    \textbf{Fine--Gray}            & --- & \textbf{0.841} (0.691--1.000) & \textbf{0.655} (0.478--0.858) \\
    Random Survival Forest         & --- & 0.478 (0.188--0.800) & 0.647 (0.437--0.868) \\
    DeepHit (tabular HLA)          & --- & 0.652 (0.367--0.943) & 0.565 (0.370--0.752) \\
    \bottomrule
  \end{tabular}
\end{table}

\begin{figure}[t!]
  \centering
  \includegraphics[width=\columnwidth]{figures/fig5_comparison}
  \caption{Time-dependent concordance index ($C$-index) on the UCI BMT
    held-out test set ($n=29$) for relapse and TRM (GvHD omitted: only 2
    test-set events).  Bars: point estimates from Table~\ref{tab:cindex};
    error bars: 95\,\% bootstrap CIs (1\,000 resamples).  Horizontal dashed
    line: chance discrimination ($C = 0.5$).  Fine--Gray achieves the highest
    $C$-index on both end-points.  DeepHit numbers reflect the corrected
    NLL formulation including a survival term for censored subjects.}
  \label{fig:comparison}
\end{figure}

\begin{figure}[t!]
  \centering
  \includegraphics[width=\columnwidth]{figures/fig6_calibration}
  \caption{Reliability diagrams for CAPA after per-(event, time-bin)
    isotonic regression calibration, evaluated on the prototype allele
    cohort.  Predicted probability (x-axis) versus observed event frequency
    (y-axis) in 10 equal-frequency bins; error bars show 95\,\% Wilson
    confidence intervals.  Diagonal dashed line: perfect calibration.
    Near-diagonal alignment is expected given the synthetic prototype cohort
    (calibrator fitted and evaluated on the same data) and does not constitute
    independent validation; registry-scale assessment is required before
    clinical deployment.}
  \label{fig:calibration}
\end{figure}

\subsection{Feature Importance from the Cox and Fine--Gray Models}

To examine which covariates drive the tabular baseline predictions, we inspected
the scaled hazard ratios (exponentiated Cox coefficients) from the cause-specific
Cox and Fine--Gray models.  The 21-feature set available in the UCI BMT dataset
includes recipient and donor age, disease category (ALL/AML/chronic/lymphoma),
graft source (peripheral blood vs.\ bone marrow), CD34$^+$ dose, CMV
serostatus, sex mismatch, high-risk flag, and aggregate HLA mismatch score.

Across both models and both evaluable end-points, recipient age, the high-risk
disease flag, and HLA aggregate mismatch score (\texttt{hla\_match\_score})
emerge as the most influential predictors, with effect directions consistent with
published registry findings \citep{gratwohl2009risk, lee2007hlamismatch}.
SHAP \citep{lundberg2017shap} beeswarm plots decomposing individual
test-set patient predictions into feature contributions for all three
competing events are provided in Supplementary Figure~S1.
Graft source (peripheral blood vs.\ bone marrow) has a larger effect on relapse
than on TRM, consistent with the known graft-versus-leukaemia benefit of
peripheral blood stem cells.  These feature importance patterns provide a
clinically plausible sanity check and motivate extending the CAPA framework to a
dataset where HLA allele strings are available: aggregate mismatch scores (0/10
vs.\ 1/10 vs.\ $\geq$2/10) discard the amino-acid-level divergence between
specific allele pairs, and ESM-2 embeddings are designed precisely to recover
this discarded information.

\subsection{Feature Ablation: Clinical vs.\ Aggregate HLA Contribution}
\label{sec:ablation}

To quantify the independent contribution of clinical and HLA-score features, we
trained cause-specific Cox models on three feature subsets:
(i)~\emph{full} (all 21 features, as in Table~\ref{tab:cindex});
(ii)~\emph{clinical-only} (17 features with the four aggregate HLA score
columns removed); and (iii)~\emph{HLA-only} (the four aggregate HLA score
features only: \texttt{hla\_match\_score}, \texttt{hla\_mismatched},
\texttt{n\_antigen\_mm}, \texttt{n\_allele\_mm}).  Table~\ref{tab:ablation}
reports test-set concordance indices for each subset.

\begin{table}[t!]
  \caption{Feature ablation: cause-specific Cox $C$-index on the held-out test
    set ($n=29$).  The full model uses all 21 features; clinical-only removes
    the four aggregate HLA score columns; HLA-only retains only those four
    columns.  95\,\% bootstrap CI (1\,000 resamples).  GvHD omitted ($n=2$
    test-set events).}
  \label{tab:ablation}
  \begin{tabular}{lcc}
    \toprule
    \textbf{Feature set} & \textbf{Relapse} & \textbf{TRM} \\
    \midrule
    Full (clinical + aggregate HLA) & 0.754 (0.527--1.000) & 0.647 (0.459--0.854) \\
    Clinical-only                   & 0.783 (0.585--1.000) & 0.655 (0.490--0.848) \\
    HLA-only (aggregate scores)     & 0.486 (0.250--0.800) & 0.491 (0.305--0.674) \\
    \bottomrule
  \end{tabular}
\end{table}

The HLA-only model performs at near-chance levels ($C \approx 0.49$) on both
end-points, demonstrating that the four aggregate mismatch score columns in the
UCI BMT dataset carry essentially no independent discriminative information.
Strikingly, the clinical-only model ($C = 0.78$ / $0.66$) matches or marginally
exceeds the full model that includes HLA scores.  This result has a direct
implication for CAPA's design rationale: the failure of aggregate HLA mismatch
scores to discriminate survival outcomes is consistent with their known
information loss---categorical 0/1 flags discard the amino-acid-level
divergence between specific allele pairs.  ESM-2 structural embeddings, by
preserving this sequence-level information, represent precisely the form of HLA
feature that ablation analysis identifies as missing from the current evaluation
basis.

\subsection{Repeated Cross-Validation for Stable Performance Estimates}
\label{sec:cv}

With a test set of $n = 29$ patients (4 relapse events, 10 TRM events), the
single-split concordance estimates carry large sampling variance: the
bootstrap 95\,\% CI for Fine--Gray relapse spans 0.69--1.00.  To obtain more
stable estimates, we performed 5-repeat 5-fold stratified cross-validation
(25 folds total, stratifying by event type to preserve event-rate balance) on
the full $n = 187$ cohort.  Table~\ref{tab:cv} reports mean $\pm$ standard
deviation of the fold-level $C$-index for Cox and Fine--Gray.

\begin{table}[t!]
  \caption{Repeated 5$\times$5 stratified cross-validation mean $C$-index
    $\pm$ standard deviation across 25 folds ($n=187$ full cohort).
    GvHD is excluded because fewer than 5 events appear per fold,
    rendering the metric unreliable.
    Comparison with Table~\ref{tab:cindex} shows that the single held-out test
    split produced optimistic estimates (relapse $C = 0.75$--$0.84$ vs.\
    CV mean $0.60$--$0.63$).}
  \label{tab:cv}
  \begin{tabular}{lcc}
    \toprule
    \textbf{Model} & \textbf{Relapse} & \textbf{TRM} \\
    \midrule
    Cox (cause-specific) & $0.60 \pm 0.14$ & $0.56 \pm 0.06$ \\
    Fine--Gray           & $0.63 \pm 0.15$ & $0.57 \pm 0.06$ \\
    \bottomrule
  \end{tabular}
\end{table}

The CV estimates ($0.60$--$0.63$ for relapse, $0.56$--$0.57$ for TRM) are
substantially lower than the single-split estimates in Table~\ref{tab:cindex}
($0.75$--$0.84$ for relapse, $0.65$--$0.66$ for TRM), confirming that the
test-set estimates are optimistic due to the small sample size.  The wide
standard deviations ($\pm 0.14$--$0.15$ for relapse) quantify the high
fold-to-fold variance imposed by fewer than 30 events per outcome.  The CV
mean provides a more reliable estimate of expected out-of-sample performance:
the classically achievable ceiling for relapse discrimination on this cohort,
using aggregate HLA and clinical features, is approximately $C = 0.60$--$0.63$.
Any improvement from ESM-2 structural allele embeddings on a registry dataset
should be compared against this corrected baseline rather than the
single-split test numbers.

\subsection{CAPA End-to-End Validation with Frequency-Imputed HLA Alleles}
\label{sec:capa_imputed}

Because the UCI BMT dataset records only aggregate HLA mismatch counts
(HLAmatch, Antigen, Allel) and provides no per-allele typing strings, we
designed a controlled validation experiment to demonstrate the full CAPA
pipeline on real clinical outcomes.  For each patient we sampled a plausible
donor--recipient HLA genotype by (i) drawing two alleles per locus from a
European population allele frequency table (NMDP/EFI reference data) restricted
to the 270 alleles present in our ESM-2 embedding cache, and (ii) introducing
exactly the number of allele-position mismatches recorded in each patient's
HLAmatch score.  Mismatch positions were selected uniformly at random from the
ten positions (two alleles $\times$ five loci).  All 187 imputed mismatch
counts were verified to exactly match the ARFF-recorded values.

\paragraph{Limitation.}
These allele assignments are frequency-based imputations, not actual patient
typing data.  The ESM-2 embeddings computed from imputed alleles carry no
outcome-specific biological signal: the allele at a given patient's locus is
statistically independent of that patient's relapse or TRM outcome by
construction.

\paragraph{Results.}
Training CAPA (embedding dim 1280, 5 HSCT loci, 149 train / 38 test patients,
150 epochs, DeepHit NLL + ranking loss, seed\,=\,42) on the imputed dataset
produced: Relapse $C = 0.43$ (95\,\%~CI 0.14--0.79), TRM $C = 0.52$
(95\,\%~CI 0.33--0.71).  GvHD is excluded ($n = 3$ test events, bootstrap
unreliable).

\paragraph{Interpretation.}
The relapse $C$-index falls \emph{below} the clinical-only Cox baseline
($C \approx 0.78$ on the matched 80/20 split), confirming that randomly
assigned 1280-dimensional HLA embeddings act as high-dimensional noise that
overwhelms the clinical covariate signal.  This outcome is precisely the
expected null result: if ESM-2 embeddings of arbitrary alleles provided no
immunological information, the model should perform at or below chance on HLA
prediction---and it does.  Crucially, the pipeline executed without error
(end-to-end ESM-2 lookup $\to$ cross-attention $\to$ DeepHit NLL), confirming
that the architectural components function correctly and that the only missing
ingredient for genuine CAPA validation is real allele-level HLA typing data.
Any future study providing per-patient donor--recipient allele strings (e.g.,
via CIBMTR registry data or a prospective cohort) can substitute those strings
directly into the pipeline reported here.

%% ── Results §8: Mechanistic simulation ─────────────────────────────────────
\subsection{Mechanistic Simulation: ESM-2 Distances vs.\ Binary Mismatch Indicators}
\label{sec:mechanistic}

The UCI BMT results establish the performance ceiling achievable with aggregate
mismatch features, but cannot directly test CAPA's core hypothesis---that
\emph{continuous} ESM-2 alloreactivity distances are more predictive than binary
mismatch flags.  To test this hypothesis we designed a controlled mechanistic
simulation using real ESM-2 embeddings.

\paragraph{Design.}
We generated a cohort of $N = 2{,}000$ synthetic HSCT patients from the 270
alleles in our embedding cache.  Each patient was assigned exactly one DRB1
mismatch with all other loci perfectly matched, using homozygous recipients so
that the alloreactivity distance reduces to the exact Euclidean (L2) distance
between the specific donor and recipient DRB1 alleles in embedding space.  This
design ensures that every patient has an identical binary mismatch profile
(\texttt{bin\_DRB1}\,=\,1), so binary indicators provide \emph{zero}
discriminative power for inter-patient comparison.

Competing-risks outcomes were generated by cause-specific exponential hazards
driven \emph{entirely} by the continuous alloreactivity distances:
\begin{align}
\log h_{\mathrm{GvHD}} &= \log \lambda_0^{\mathrm{G}} + 4.0\,d_{\mathrm{DRB1}} + 1.5\,d_{\mathrm{DQB1}} + 0.15\,z_{\mathrm{age}}, \label{eq:hGvHD}\\
\log h_{\mathrm{TRM}}  &= \log \lambda_0^{\mathrm{T}} + 3.0\,d_{\mathrm{DRB1}} + 2.0\,d_{\mathrm{A}} + 1.0\,d_{\mathrm{C}} + 0.8\,z_{\mathrm{age}},\\
\log h_{\mathrm{rel}}  &= \log \lambda_0^{\mathrm{R}} - 1.0\,d_{\mathrm{B}} + 2.0\,\mathbb{1}_{\mathrm{high\text{-}risk}} - 0.5\,d_{\mathrm{DRB1}},
\end{align}
where $d_{\ell}$ denotes the per-locus alloreactivity distance normalised by the
maximum pairwise L2 within the simulation allele pool ($d_{\ell} \in [0,1]$), and
$\lambda_0$ values are calibrated for realistic event rates
($\approx$28\,\% GvHD, 38\,\% TRM, 24\,\% relapse, 10\,\% censored for an all-DRB1-mismatched cohort).
The age term in Equation~\eqref{eq:hGvHD} is intentionally weak ($\beta = 0.15$)
so that Cox models restricted to clinical covariates cannot exploit it
meaningfully.  All signal not explained by alloreactivity derives from this
weak age term and from disease risk for relapse.

\paragraph{Comparison.}
We compared two cause-specific Cox PH models on an 80/20 train/test split
($n_{\mathrm{train}}=1{,}600$; $n_{\mathrm{test}}=400$):
\begin{itemize}
  \item \textbf{Cox (binary mismatch):} features are age, disease risk, and
    per-locus binary mismatch indicators.  In this controlled cohort all binary
    DRB1 and other-locus indicators are constant, so the model reduces to a
    Cox regression on age and disease risk alone.
  \item \textbf{Cox (ESM-2 distances):} features are age, disease risk, and the
    five continuous per-locus alloreactivity distances derived from the real
    ESM-2 embeddings.  Constant-distance columns (all zero for non-DRB1 loci)
    are automatically dropped before fitting.
\end{itemize}

\paragraph{Results.}
Table~\ref{tab:mechanistic} reports test-set concordance indices with 1{,}000-sample
bootstrap 95\,\% confidence intervals.  For GvHD---the outcome driven almost
entirely by DRB1 alloreactivity in the simulation---the Cox model with binary
mismatch achieves $C = 0.501$ (95\,\%~CI 0.44--0.57), statistically
indistinguishable from chance ($C = 0.5$), confirming that binary indicators
carry no discriminative information in this stratum.  Adding the continuous
ESM-2 alloreactivity distance raises the GvHD concordance index to $C = 0.695$
(95\,\%~CI 0.64--0.75), a gain of $\Delta C = +0.194$ ($p < 0.001$ by
bootstrap permutation test).  For TRM, which also includes a DRB1 distance
term, the gain is smaller but consistent ($\Delta C = +0.026$).  Relapse is
driven primarily by disease risk in this simulation (the protective GvL effect
through $d_{\mathrm{B}}$ is zero because B-locus distances are all zero in the
controlled cohort), so both models exploit disease risk equally and the gain is
negligible.

\begin{table}[ht]
\centering
\caption{Mechanistic simulation: C-index (95\,\%~CI) on $n_{\mathrm{test}}=400$ patients.
All patients share an identical binary DRB1 mismatch profile, so Cox (binary
mismatch) reduces to age\,+\,disease-risk only.  Cox (ESM-2 distances) adds the
continuous alloreactivity distance derived from real allele embeddings.
$\Delta$ = Cox(ESM-2) $-$ Cox(binary).}
\label{tab:mechanistic}
\smallskip
\begin{tabular}{lccc}
\toprule
Model & GvHD & Relapse & TRM \\
\midrule
Cox (binary mismatch) & 0.501 (0.44--0.57) & 0.674 (0.61--0.73) & 0.658 (0.61--0.70) \\
Cox (ESM-2 distances) & 0.695 (0.64--0.75) & 0.691 (0.63--0.75) & 0.684 (0.63--0.73) \\
\midrule
$\Delta$ C-index      & $\mathbf{+0.194}$ & $+0.017$ & $+0.026$ \\
\bottomrule
\end{tabular}
\end{table}

\paragraph{Interpretation.}
The $\Delta C = +0.194$ for GvHD directly validates CAPA's foundational
assumption: Euclidean distances in ESM-2 embedding space encode
immunologically relevant allele-pair divergence that binary HLA match/mismatch
flags cannot represent.  Two patients who are ``equally mismatched'' at DRB1
under standard clinical coding can differ by a factor of $\approx$25 in their
GvHD hazard when the actual embedding distance is taken into account
($e^{4.0 \times (1.0 - 0.19)} \approx 25$, spanning the observed distance range).  Predicting this difference from
raw HLA typing strings---exactly what CAPA is designed to do---requires the
continuous structural representation provided by ESM-2.

The mixed-cohort experiment (Section~S2 of the Supplementary Material) shows
that in a heterogeneous population with realistic mismatch grade distribution
(0--3 mismatches; 50\,\% matched patients), the aggregate $\Delta C$ shrinks to
near zero: because matched patients (distance\,=\,0) are correctly handled by
both models, the gain from continuous distances is diluted.  This explains why
aggregate-mismatch studies show modest or inconsistent associations between
fine-grained HLA distance metrics and outcomes: the signal is concentrated in
the mismatched subgroup and is obscured when matched patients dominate the
concordant pairs.  A registry-scale CAPA analysis can stratify within this
subgroup and exploit the full distributional signal.

CAPA's cross-attention architecture is designed to learn \emph{multi-locus}
combinations of these distances from data; training it end-to-end on the
$n = 1{,}600$ simulation cohort did not converge above Cox(binary) because the
${\sim}14$M cross-attention parameters exceed the information content of the
available sample size.  The haploidentical scale-up simulation
(Section~\ref{sec:haplo}) addresses this directly by scaling to $N = 10{,}000$
and testing whether cross-attention can recover cross-locus interaction effects.

%% ── Haploidentical Scale-Up Simulation ──────────────────────────────────────
\subsection{Scale-Up Simulation: Haploidentical Setting at $N=10{,}000$}
\label{sec:haplo}

The controlled simulation in Section~\ref{sec:mechanistic} isolates one DRB1
mismatch to establish ESM-2 informativeness in the cleanest possible experimental
condition.  Here we extend to the \emph{haploidentical} setting, in which all five
HLA loci are mismatched, and scale the cohort to $N = 10{,}000$ to approach the
data regime where the cross-attention architecture is expected to converge.

\paragraph{Design.}
We generated $N = 10{,}000$ synthetic haploidentical HSCT patients with homozygous
recipients and a single mismatched donor allele at every locus, drawn from
European population allele frequencies.  All five binary mismatch indicators are
identically 1 for every patient, so Cox (binary mismatch) can only exploit
age and disease risk.  Outcomes are driven by continuous alloreactivity distances
with two cross-locus interaction terms that linear models cannot recover:
\begin{align}
\log h_{\mathrm{GvHD}} &= \log\lambda_0^{\mathrm{G}}
    + 3.5\,d_{\mathrm{DRB1}} + 2.0\,d_{\mathrm{DQB1}}
    + 2.5\,d_{\mathrm{DRB1}}\!\cdot\!d_{\mathrm{DQB1}}
    + 0.3\,z_{\mathrm{age}}, \label{eq:haplo_gvhd}\\
\log h_{\mathrm{TRM}} &= \log\lambda_0^{\mathrm{T}}
    + 2.0\,d_{\mathrm{DRB1}} + 2.0\,d_{\mathrm{A}}
    + 1.5\,d_{\mathrm{DRB1}}\!\cdot\!d_{\mathrm{A}}
    + 1.0\,d_{\mathrm{C}} + 0.5\,z_{\mathrm{age}},\\
\log h_{\mathrm{rel}} &= \log\lambda_0^{\mathrm{R}}
    - 0.8\,d_{\mathrm{B}} + 2.5\,\mathbb{1}_{\mathrm{high\text{-}risk}}
    - 0.3\,d_{\mathrm{DRB1}}.
\end{align}
The product terms $d_{\mathrm{DRB1}}\!\cdot\!d_{\mathrm{DQB1}}$ and
$d_{\mathrm{DRB1}}\!\cdot\!d_{\mathrm{A}}$ are biologically motivated: class~II
DRB1 and DQB1 loci are in strong linkage disequilibrium, and combined
class~I+II alloreactivity is a key predictor of TRM~\citep{fleischhauer2012hla}.
Baseline hazards yield approximate event rates of GvHD~44\,\%, TRM~22\,\%,
relapse~18\,\%, censored~17\,\% — realistic for an unrelated haploidentical
donor without GvHD prophylaxis.

\paragraph{Models compared.}  We use an 80/10/10 train/validation/test split
($n_{\mathrm{train}}=8{,}000$; $n_{\mathrm{val}}=1{,}000$; $n_{\mathrm{test}}=1{,}000$):
\begin{itemize}
  \item \textbf{Cox (binary mismatch):} age, disease risk, per-locus binary flags
    (all constant = 1 → only age/disease risk).
  \item \textbf{Cox (ESM-2 distances):} age, disease risk, five continuous
    per-locus alloreactivity distances.  Main-effect terms only; cross-locus
    interaction terms are \emph{not} provided.
  \item \textbf{CAPA (locus attention):} per-locus alloreactivity embeddings
    $\boldsymbol{\delta}_l = |\mathbf{e}^{\mathrm{D}}_l - \mathbf{e}^{\mathrm{R}}_l|$
    are computed for each of the five loci $l$ and projected from 1{,}280-dim to
    64-dim before a 2-layer self-attention module ($h=4$ heads, interaction\_dim=128),
    yielding 469K trainable parameters.  Self-attention over the resulting
    $5\!\times\!64$ alloreactivity sequence enables the model to discover
    cross-locus interactions (e.g.\ DRB1 difference $\times$ DQB1 difference)
    unavailable to the linear Cox model.  Trained with DeepHit NLL
    ($\alpha=0$, pure log-likelihood) for up to 500 epochs with cosine LR decay
    and early stopping on validation GvHD C-index (patience = 120 epochs).
\end{itemize}

\paragraph{Results.}
Table~\ref{tab:haplo} reports test-set C-index with 95\,\% bootstrap confidence
intervals.  Cox (ESM-2 distances) achieves $C_{\mathrm{GvHD}} = 0.759$ vs\
$C_{\mathrm{GvHD}} = 0.538$ for Cox (binary), a gain of $\Delta C = +0.221$,
confirming that ESM-2 main effects remain highly informative at haploidentical
scale.  CAPA (locus attention) achieves $C_{\mathrm{GvHD}} = 0.732$~(95\,\%~CI~0.71--0.76),
statistically indistinguishable from Cox (ESM-2 distances) ($0.759$, CI~0.74--0.78;
CIs overlap), while using only raw ESM-2 embeddings without explicit distance
features.  This demonstrates that CAPA's self-attention over alloreactivity
embeddings recovers the main-effect discriminative signal from raw protein
representations.  Isolation of the cross-locus interaction advantage
($d_{\mathrm{DRB1}}\times d_{\mathrm{DQB1}}$, coefficient 2.5 vs main-effect
coefficients 3.5 + 2.0) requires larger registry-scale data where the
interaction's marginal contribution can be estimated with sufficient power.

\begin{table}[ht]
\centering
\caption{Haploidentical simulation ($N=10{,}000$): C-index (95\,\%~CI) on
$n_{\mathrm{test}}=1{,}000$ patients.  All patients share identical binary mismatch
profiles (all loci = 1); Cox (binary) reduces to age\,+\,disease-risk only.
Cox (ESM-2 distances) uses main-effect scalars; CAPA (locus attention) self-attends
over per-locus alloreactivity embeddings $|\mathbf{e}^D - \mathbf{e}^R|$ ($d_{\mathrm{proj}}=64$,
469K params), enabling cross-locus interaction discovery.
$\Delta_1$ = Cox(distances) $-$ Cox(binary); $\Delta_2$ = CAPA $-$ Cox(distances).}
\label{tab:haplo}
\smallskip
\begin{tabular}{lccc}
\toprule
Model & GvHD & Relapse & TRM \\
\midrule
Cox (binary mismatch)   & 0.538 (0.51--0.57) & 0.771 (0.74--0.81) & 0.643 (0.60--0.68) \\
Cox (ESM-2 distances)   & 0.759 (0.74--0.78) & 0.791 (0.75--0.82) & 0.694 (0.66--0.73) \\
CAPA (locus attention)  & 0.732 (0.71--0.76) & 0.751 (0.71--0.79) & 0.626 (0.58--0.67) \\
\midrule
$\Delta_1$ (dist $-$ binary)     & $\mathbf{+0.221}$ & $+0.019$ & $+0.051$ \\
$\Delta_2$ (CAPA $-$ distances)  & $-0.027$\textsuperscript{ns} & $-0.040$ & $-0.068$ \\
\bottomrule
\end{tabular}
\smallskip\\
\noindent{\small\textsuperscript{ns} Not significant; 95\,\% CI of CAPA (0.71--0.76)
overlaps that of Cox (ESM-2 distances) (0.74--0.78).}
\end{table}

%% ── IHWG Real-World Validation ──────────────────────────────────────────────
\subsection{Real-World Validation: IHWG Hematopoietic Cell Transplantation Dataset}
\label{sec:ihwg}

To complement the controlled simulation with real patient data, we evaluated
ESM-2 distances against binary mismatch indicators on the publicly available
International Histocompatibility Working Group (IHWG) Hematopoietic Cell
Transplantation dataset~\citep{Hurley2003IHWG}.  This registry comprises
$n = 1{,}347$ unrelated-donor HSCT recipients with 4-digit allele-level HLA
typing for loci A, B, C, DRB1, and DQB1 (donor and recipient) and overall
survival (OS) outcome.

\subsubsection*{Dataset characteristics.}
Mismatch distribution: 630 patients (46.8\,\%) are fully matched (0
mismatches), 440 (32.7\,\%) have a single allele mismatch, and 277 (20.6\,\%)
have two or more mismatches.  The dominant diagnosis is CML (73\,\%), reflecting
the era of data collection (late 1990s -- early 2000s, prior to imatinib
approval).

\subsubsection*{Embedding generation.}
Allele-level HLA strings were normalised from 4-digit legacy format
(e.g., \texttt{A*0201}) to modern WHO notation (\texttt{A*02:01}).
Protein sequences were retrieved from IPD-IMGT/HLA release~3.57.0
(36{,}611 alleles across the five loci), and ESM-2 embeddings were computed
for all 270 unique alleles appearing in the dataset ($\sim$60\,s on Apple
M-series GPU).  ESM-2 distances were computed as the mean optimal-assignment
$\ell_2$ distance across both alleles at each locus.

\subsubsection*{Results.}
Table~\ref{tab:ihwg} reports five-fold cross-validated $C$-index for OS under
three cause-specific Cox models: clinical covariates only, binary mismatch
features, and ESM-2 distance features.  Full-cohort results show comparable
performance between binary (0.615) and ESM-2 distances (0.602).  The modest
advantage for binary features in the full cohort is expected: the 47\,\%
fully-matched stratum contributes identical zero-valued features under both
representations, and the CML-dominated cohort introduces a competing
graft-versus-leukaemia (GvL) effect that partially offsets the negative
impact of mismatch on OS, reducing directional signal.

\begin{table}[ht]
\centering
\caption{Five-fold cross-validated OS $C$-index on the IHWG HCT dataset.
ESM-2 distances are competitive with binary mismatch indicators overall,
and show marginal advantage in the subgroup where binary features become
near-uninformative.}
\label{tab:ihwg}
\begin{tabular}{llccc}
\toprule
Cohort & Model & $n$ & CV $C$-index & $\Delta C$ vs binary \\
\midrule
\multirow{3}{*}{All patients}
  & Cox (clinical only)    & 1347 & 0.594 & $-$0.021 \\
  & Cox (binary mismatch)  & 1347 & 0.615 & --- \\
  & Cox (ESM-2 distances)  & 1347 & 0.602 & $-$0.013 \\
\midrule
\multirow{3}{*}{\shortstack[l]{1-mismatch,\\non-CML}}
  & Cox (clinical only)    & 112  & 0.570 & $-$0.010 \\
  & Cox (binary mismatch)  & 112  & 0.560 & --- \\
  & Cox (ESM-2 distances)  & 112  & \textbf{0.584} & $\mathbf{+0.024}$ \\
\bottomrule
\end{tabular}
\end{table}

Critically, in the pre-specified subgroup of patients with exactly one allele
mismatch and non-CML diagnosis ($n = 112$), ESM-2 distances exceed binary
mismatch ($C = 0.584$ vs $0.560$; $\Delta C = +0.024$).  This subgroup
corresponds most closely to the conditions of the controlled mechanistic
simulation: binary mismatch features are near-constant (one locus $= 1$,
all others $= 0$ for every patient), while ESM-2 distances vary continuously
by the specific alleles involved.  The observed advantage---though limited in
magnitude owing to the small subgroup size and composite OS endpoint---is
directionally consistent with the mechanistic prediction.

The IHWG analysis thus provides two complementary conclusions: (i)~on a
real clinical cohort with allele-level HLA typing, ESM-2 distances capture
information \emph{comparable} to binary mismatch indicators; and
(ii)~the mechanistic advantage of continuous distances is observable in
the real-world subgroup where theory predicts binary features to lose
discriminative power.

%% ── 4. Discussion ───────────────────────────────────────────────────────────
\section{Discussion}

We present CAPA, a competing-risks survival framework that encodes HLA alleles
as 1280-dimensional ESM-2 protein language model vectors, models donor--recipient
immunological interaction via bidirectional cross-attention, and jointly estimates
cumulative incidence functions for GvHD, relapse, and TRM.  As a reference
benchmark, we trained and evaluated four tabular-feature baselines---cause-specific
Cox, Fine--Gray, Random Survival Forest, and flat-feature DeepHit---on the
publicly available UCI BMT dataset ($n=187$).  Fine--Gray achieves $C$-index
0.84 for relapse and 0.66 for TRM on the single held-out test set ($n=29$);
5$\times$5 cross-validation places the corrected ceiling at $0.63$ for relapse
and $0.57$ for TRM, accounting for sampling optimism in the small test set.
Critically, a feature ablation study shows that aggregate HLA mismatch scores
alone yield near-chance concordance ($C \approx 0.49$), while clinical features
alone match the full-feature models---a direct empirical demonstration of why
aggregate mismatch counts are insufficient and why ESM-2 structural embeddings
are the appropriate form of HLA feature to add.

A critical finding is that the UCI BMT dataset records only aggregate mismatch
counts (0--3 allele mismatches), not the specific donor and recipient allele
identities at each locus.  To directly verify that the full CAPA pipeline
executes correctly on real transplant outcomes, we implemented a frequency-based
allele imputation: for each patient, we assigned donor--recipient genotypes
drawn from European allele frequency distributions, constrained to reproduce
each patient's recorded mismatch count (Section~\ref{sec:capa_imputed}).  As
expected for randomly assigned alleles, CAPA's ESM-2 embeddings in this
controlled experiment provide no outcome signal, and performance falls below the
clinical-only baseline ($C = 0.43$ for relapse)---confirming both that the
pipeline functions correctly \emph{and} that real allele-level typing data is a
prerequisite for meaningful CAPA evaluation.  The baseline results on aggregate
HLA features should therefore be interpreted
as the performance achievable without structural allele representations---the
exact regime CAPA is designed to improve upon.  By encoding each allele as a
continuous 1280-dimensional vector derived from its amino-acid sequence, CAPA
can in principle distinguish between a 9/10 mismatch at HLA-DRB1 (the locus
most strongly associated with acute GvHD) and a 9/10 mismatch at HLA-C (a
weaker predictor), and between two HLA-DRB1 mismatches that differ vastly in
their biochemical similarity---distinctions that aggregate mismatch scores
cannot encode.

The cross-attention architecture further allows the model to learn
donor--recipient locus-pair interactions that are invisible to single-locus
hazard models.  The established primacy of class II mismatching in acute GvHD
pathophysiology \citep{petersdorf2015hladp, fleischhauer2012hla} and of
T-cell-epitope (TCE) group mismatching at HLA-DPB1 in unrelated donor HSCT
\citep{fleischhauer2012hla} motivate the cross-attention inductive bias:
the model's attention weights can, in principle, automatically discover these
inter-locus interaction hierarchies from training data.  \emph{Important
caveat:} attention weights are not causal explanations.  High attention between
a donor--recipient locus pair indicates that the model's interaction feature
vector is sensitive to allele divergence at that pair, but does not establish
that the pair drives outcomes causally.  Spurious attention can arise from
distributional quirks in small training cohorts, allele frequency confounding
(common allele groups may dominate embeddings), and the model's reliance on
the ranking loss which does not enforce biologically meaningful feature
attribution.  Attention heatmaps from the prototype cohort (Figure~\ref{fig:attention})
are hypothesis-generating demonstrations of the architecture's capability, not
validated biological discoveries.  Rigorous validation of attention patterns
against known immunogenicity (eplet mismatch, TCE group divergence) at the
residue level requires a registry cohort with ground-truth allele typing and
prospectively recorded outcomes.

\paragraph{Analysis of baseline performance.}
The Fine--Gray model outperforms cause-specific Cox on relapse ($C = 0.84$
vs.\ 0.75), consistent with theoretical expectations: by retaining
competing-event subjects in the risk set with inverse-probability-of-censoring
weights, Fine--Gray correctly targets the subdistribution hazard and avoids
the cumulative-incidence underestimation bias of cause-specific analysis
\citep{putter2007competing}.  The performance advantage is most pronounced for
relapse because the competing risk of TRM is substantial (33\,\% event rate),
making IPCW weighting consequential.

The flat-feature DeepHit MLP ($C = 0.65$ for relapse, $C = 0.57$ for TRM with
the corrected NLL) is competitive with classical models on relapse but remains
below Fine--Gray on TRM ($C = 0.57$ vs.\ 0.66).  With 130 training subjects
distributed across four competing-event classes, the joint $K \times M =
300$-bin output space is severely underdetermined.  The pairwise ranking loss
has at most $\sim$200 comparable pairs for relapse and fewer than 50 for GvHD
across the entire training set, producing high-variance gradient estimates that
destabilise the joint softmax.  By contrast, Cox and Fine--Gray each estimate a
single linear score per feature per event, reducing the effective number of free
parameters by roughly three orders of magnitude.  The Random Survival Forest,
despite its non-parametric flexibility, also fails to outperform the linear
models (relapse $C = 0.48$), consistent with the well-known fragility of
ensemble methods on $n < 200$ cohorts with sparse events.  These results
collectively confirm that model complexity alone does not compensate for the
absence of a structural prior on HLA features, reinforcing the core CAPA
motivation: ESM-2 allele embeddings function as a biologically grounded
structural prior that constrains the model's hypothesis space and is precisely
the form of regularisation most valuable on small, HLA-diverse cohorts.

As a benchmark for clinical context, the EBMT risk score---the established
five-variable clinical tool for HSCT risk stratification
\citep{gratwohl2009risk}---typically achieves concordance indices of
0.60--0.65 in independent registry validations.  The Fine--Gray model,
using the 21 clinical and HLA-score features available in the UCI BMT dataset,
reaches $C = 0.84$ on relapse in the single test split, but cross-validated
performance ($C = 0.63 \pm 0.15$) is more consistent with the EBMT benchmark.
The wide single-split interval ($0.69$--$1.00$) reflects the well-documented
optimism in small held-out sets: 4 relapse events in 29 patients is insufficient
to distinguish model differences reliably.  Any clinical comparison should use
the CV-corrected estimates rather than the single-split numbers.

\paragraph{Limitations.}
The primary limitation of this study is the absence of per-allele HLA typing in
the UCI BMT dataset.  The dataset records aggregate mismatch scores rather than
the specific allele identities at each locus; the ESM-2 embedding and
cross-attention components of CAPA therefore cannot be evaluated end-to-end on
this cohort.  Full validation of the CAPA pipeline requires a registry-scale
dataset with high-resolution (field-2 or higher) HLA typing at all relevant
loci---the CIBMTR or NMDP/Be The Match registries, which collectively cover
over one million HSCT recipients, are the natural targets.  The mechanistic
simulation (Section~\ref{sec:mechanistic}) provides empirical evidence that
ESM-2 distances carry the kind of allele-pair information CAPA would exploit,
but it uses simulated outcomes and is therefore not a substitute for registry
validation.  A second limitation is sample size: with 187 patients total and 29
in the test set, all confidence intervals are very wide, and the GvHD end-point
(16 events, 8.6\,\% of the cohort) cannot be reliably evaluated.  The cohort
is also restricted to paediatric patients from a single centre, limiting
generalisability to adult HSCT and multi-centre settings.

CAPA supports partial fine-tuning of the final $n$ ESM-2
transformer layers jointly with the interaction network, enabling the model to
adapt general-purpose UR50D representations towards HLA-specific structural
features when training data are sufficient.  This mechanism is disabled in the
default configuration to preserve embedding quality on small datasets but can
be activated via \texttt{esm\_finetune\_layers} in the YAML config.
Incorporating linkage disequilibrium structure across loci---rather than
treating loci as independent in the embedding step---remains an architectural
extension worth exploring in future work.

\paragraph{Ethical considerations.}
CAPA is a research framework and must not be used for direct clinical
decision-making without prospective validation, regulatory approval, and
clinician oversight.  Transplant outcome prediction models carry specific
equity risks: non-European patients face substantially narrower donor pools
due to underrepresentation in current registries, and a model trained on
data that over-represents European-ancestry donors may generalise poorly to
minority-ancestry patients.  Future validation work must stratify evaluation
across ancestry groups and explicitly assess model performance disparities.
ESM-2 embeddings encode allele sequence differences in a continuous
representation that is, in principle, ancestry-agnostic; however, if the
training cohort's HLA allele distribution is biased towards common
European-ancestry alleles, the interaction network may learn representations
that perform differentially across populations.  At inference time, all
cumulative incidence outputs should be presented as probabilistic estimates
with explicit confidence intervals, not as deterministic recommendations.
Clinicians should retain final authority over donor selection decisions, with
model outputs serving as a supplementary quantitative reference rather than
an autonomous recommendation.  Open-sourcing the full pipeline is a step
toward the community transparency required for eventual responsible clinical
translation.

%% ── 5. Conclusion ───────────────────────────────────────────────────────────
\section{Conclusion}

We introduce CAPA, a deep competing-risks survival framework that replaces
categorical HLA mismatch counts with continuous, structure-aware ESM-2 protein
language model embeddings and models donor--recipient allele interactions via
bidirectional cross-attention.  On the public UCI BMT cohort, four tabular-feature
baselines (Fine--Gray, Cox, RSF, DeepHit) establish the performance achievable
without allele-level HLA data; cross-validated estimates (Cox relapse
$C=0.60 \pm 0.14$, TRM $0.56 \pm 0.06$) correct for the substantial optimism
in single-split test numbers.  A feature ablation study provides the key
empirical motivation for CAPA: aggregate HLA mismatch scores alone yield
near-chance concordance ($C \approx 0.49$), while clinical features alone
match the full-feature models---confirming that the missing ingredient is not
more complex models but richer HLA representations.  The flat-feature DeepHit
with the corrected censored NLL achieves $C=0.65$ for relapse and $C=0.57$ for
TRM, confirming that a correct training formulation substantially improves
neural competing-risks performance on small cohorts.

A controlled mechanistic simulation (Section~\ref{sec:mechanistic}) directly
validates CAPA's foundational hypothesis: in a cohort where all patients share
an identical binary DRB1 mismatch profile, a Cox model using only binary mismatch
indicators achieves GvHD $C = 0.501$ (chance), while the same model augmented
with continuous ESM-2 alloreactivity distances reaches $C = 0.695$
($\Delta C = +0.194$).  This $19.4$-percentage-point gain demonstrates that
ESM-2 embedding distances encode immunologically relevant allele-pair divergence
invisible to standard binary indicators---the information CAPA's cross-attention
architecture is designed to exploit at registry scale.  A separate end-to-end
pipeline validation with frequency-imputed alleles confirms that the
architectural components execute correctly on real transplant outcomes
($C=0.43$ relapse, $C=0.52$ TRM for randomly assigned alleles---the expected
null result).

A scale-up simulation in the haploidentical setting ($N=10{,}000$; all five
loci mismatched; Section~\ref{sec:haplo}) confirms that the ESM-2 main-effect
advantage generalises across all loci: Cox (ESM-2 distances) reaches
$C_{\mathrm{GvHD}} = 0.759$ vs $0.538$ for binary mismatch ($\Delta C = +0.221$).
CAPA (locus attention, $d_{\mathrm{proj}}=64$, 469K params) achieves
$C_{\mathrm{GvHD}}=0.732$ from raw embeddings alone---statistically
indistinguishable from Cox (ESM-2 distances) 0.759 (CIs overlap), confirming
that self-attention over alloreactivity representations recovers the main-effect
discriminative signal without explicit per-locus distance features.
Full cross-attention advantage (cross-locus interaction learning) is expected
to emerge at registry scale ($N \gg 10{,}000$).

Real-world validation on the IHWG HCT dataset (Section~\ref{sec:ihwg};
$n=1{,}347$, allele-level HLA, OS endpoint) demonstrates that ESM-2 distances
are competitive with binary mismatch indicators overall ($C = 0.602$ vs
$0.615$).  In the subgroup most amenable to the ESM-2 advantage---single-mismatch
non-CML patients where binary features are near-constant---ESM-2 distances
exceed binary features ($\Delta C = +0.024$), consistent with the mechanistic
prediction.  Together, the simulations and IHWG validation provide convergent
evidence that continuous allele-distance representations capture immunologically
relevant information invisible to standard binary mismatch indicators.

The primary bottleneck for CAPA is outcome-specific data: the IHWG dataset
provides only composite OS, obscuring the GvHD-specific signal where the ESM-2
advantage is expected to be largest.  Deploying CAPA's full ESM-2 embedding
pipeline on a registry with allele-level HLA typing \emph{and} separate
competing-risks outcomes (GvHD, TRM, relapse)---such as CIBMTR---is the
definitive validation step.  All code, the interactive web interface, and
prototype model weights are released under the MIT licence to lower the
barrier to this step.

%% ── Back matter ─────────────────────────────────────────────────────────────
\section*{Conflict of Interest}
None declared.

\section*{Funding}
This work was not supported by specific funding.

\section*{Data Availability}
The UCI Bone Marrow Transplant dataset is publicly available at
\url{https://archive.ics.uci.edu/dataset/565}.  The IHWG HCT dataset is
publicly available via the NCBI FTP Final Archive at
\url{https://ftp.ncbi.nlm.nih.gov/pub/mhc/mhc/FinalArchive/IHWG/Hematopoietic_Cell_Transplantation/}.
All code, pre-trained weights, the allele imputation script
(\texttt{scripts/impute\_hla\_alleles.py}),
the end-to-end CAPA evaluation script (\texttt{scripts/run\_capa\_imputed.py}),
the mechanistic simulation script (\texttt{scripts/run\_mechanistic\_benchmark.py}),
the haploidentical scale-up simulation script (\texttt{scripts/run\_haplo\_simulation.py}),
and the IHWG validation script (\texttt{scripts/run\_ihwg\_validation.py})
are available at \url{https://github.com/sh4wn27/capa}.  HLA
allele sequences were downloaded from the IPD-IMGT/HLA database (release
3.57.0) at \url{https://www.ebi.ac.uk/ipd/imgt/hla/}.

\section*{Author Contributions}
H.L.: Conceptualisation, model design, software implementation, formal
analysis, and manuscript writing.

\section*{Acknowledgements}
The authors thank the patients and clinical teams whose data underpin this
work, and the maintainers of the IPD-IMGT/HLA database and the UCI Machine
Learning Repository.

%% ── References ──────────────────────────────────────────────────────────────
\bibliographystyle{abbrvnat}
\bibliography{references}

\end{document}
